% Canonical node 019 · P003; current section path 1.3.4.1.
% Proof owner: M001
\cnnaProofNodeHeading{019}{P003}{Birth-cut canonicity and partition well-formedness}
\cnnaNodePlacement{1.3.4.1}{D2}
Let $X_n$ be a C005 state, let $s_{n+1}$ be its C004 successor, and let
$\mathcal C_n(s_{n+1})=(B_n,I_n)$ be the M001 cut.  P003 certifies that this
construction is a genuine ordered partition of the already-born carrier:
\[
 B_n\cap I_n=\varnothing,
 \qquad
 B_n\cup I_n=V_n,
 \qquad
 c_{n+1}\notin B_n\cup I_n.
\]
Every entry of both lists is born.  The two lists preserve C005/C018 order
because M001 defines them by filtering the canonical carrier rather than by a
new sorting operation.  The predicate \emph{IsCanonicalCut} identifies exactly
the two filtered lists, so any two canonical witnesses are equal.

The certificate is implemented as internal theorems of the M001 owner module,
not as a second mathematical construction or a parallel proof DAG.  The
relevant theorem chain is
\[
\texttt{boundary\_node\_born},\ 
\texttt{interior\_node\_born},\ 
\texttt{boundary\_interior\_disjoint},\ 
\texttt{carrier\_covered},\ 
\texttt{child\_not\_in\_boundary},\ 
\texttt{child\_not\_in\_interior},\ 
\texttt{unique}.
\]
C004 directly supplies the verified least-open successor, while C018 supplies
the order needed by M001.  P003 therefore closes the M001 owner gate without a
P002 dependency.
